% Image credits for the photographs and AI illustrations of Book 2
% (University Chemistry, Year 1). Machine-readable license records live in
% images/book2/CREDITS.md; keep the two files in sync. The Book 2 writer
% owns this file: add one \item per photograph (in an itemize, after the
% first paragraph) as photographs land.
\chapter*{图片致谢}

本书中的插图均为原创。照片（如有使用）依其各自的免费许可协议使用，谨向其作者致谢；每张照片的完整来源链接和许可协议网址，均列于本书代码库中的 \texttt{images/book2/CREDITS.md}。
危险象形图取自全球化学品统一分类和标签制度（Globally Harmonized System of
Classification and Labelling of Chemicals），其绘制形式与 \TeX\ Live 的 \texttt{ghsystem} 宏包一致。

\bigskip
除照片和原创插图外，部分图形为人工智能生成的插图，由 OpenAI 图像生成功能依据本书编辑撰写的提示词制作，并在收入前经过化学准确性审核。每幅生成图像的完整提示词均记录于代码库中的 \texttt{images/book2/ai/PROMPTS.md}。

\bigskip
\noindent 照片，按章列出：
\begin{itemize}\small
  \item 第~1 章：尼尔斯·玻尔（Niels Bohr），约 1922 年，AB Lagrelius \& Westphal，公有领域（维基共享资源）。
  \item 第~2 章：莱纳斯·鲍林（Linus Pauling），1962 年，Nobel Foundation，公有领域（维基共享资源）。
  \item 第~5 章：天然铜样本，作者 Flipin，OG，CC0（维基共享资源）。
  \item 第~6 章：岩盐晶体、萤石和八面体金刚石，作者 James St. John，CC BY 2.0；石墨，作者 Zbynek Burival，CC BY-SA 4.0；雪晶，作者 Wilson A. Bentley（1902），公有领域（均为维基共享资源）。
  \item 第~8 章：斯范特·阿伦尼乌斯（Svante Arrhenius），1909 年，Meisenbach Riffarth \& Co.，公有领域（维基共享资源）。
  \item 第~11 章：氯化银、溴化银和碘化银沉淀，作者 Rrausch1974，CC BY-SA 3.0（维基共享资源）。
  \item 第~12 章：氢氧化铜(II) 悬浊液，作者 Alvy16，CC BY 4.0；铜氨溶液，作者 Chemicalinterest，公有领域（均为维基共享资源）。
  \item 第~13 章：瓦尔特·能斯特（Walther Nernst），1889 年，Smithsonian Libraries，公有领域（维基共享资源）。
  \item 第~16 章：雅各布斯·亨德里克·范托夫（Jacobus Henricus van 't Hoff），照片发表于 1911 年，公有领域（维基共享资源）。
  \item 第~21 章：维克多·格利雅（Victor Grignard），公有领域（维基共享资源）。
  \item 第~25 章：弗拉基米尔·马尔科夫尼科夫（Vladimir Markovnikov），肖像取自 1905 年发表的一则讣告，公有领域（维基共享资源）。
  \item 第~26 章：阿塔卡马盐沼的锂盐池，NASA Earth Observatory，公有领域；矿物油中的钠，作者 Justin Urgitis，CC BY-SA 2.5；火焰颜色，作者 Hegelrast，CC BY-SA 4.0（均为维基共享资源）。
  \item 第~27 章：豆状铝土矿，作者 James St. John，CC BY 2.0；硼砂，作者 Aram Dulyan，公有领域；石英，作者 Stephanie Clifford，CC BY 2.0；白磷，Hi-Res Images of Chemical Elements，CC BY 3.0；弗里茨·哈伯（Fritz Haber），The Nobel Foundation，公有领域（均为维基共享资源）。
  \item 第~28 章：卡瓦伊真火山口的硫矿开采，作者 S\'emhur，CC BY-SA 4.0；安瓿瓶中的氯，作者 W. Oelen，CC BY-SA 3.0；安瓿瓶中的溴，作者 Jurii，CC BY 3.0；碘晶体，作者 Dnn87，CC BY 3.0（均为维基共享资源）。
  \item 第~29 章：加热台（科弗勒台），作者 HBR，CC BY-SA 3.0；阿贝折射仪，作者 Foreade，CC BY-SA 4.0（均为维基共享资源）。
\end{itemize}
\clearpage
